% !TEX root = ../main.tex
\begin{ouverture}[equilibre]{7}{Équilibre chimique et sens d'évolution}{Partie 1 · Constitution et transformations de la matière}
{Taux d'avancement final\enspace·\enspace Équilibre dynamique\enspace·\enspace Quotient de réaction\enspace·\enspace Critère d'évolution}
\objectif{Relier le caractère non total d'une transformation à la présence de tous les réactifs et produits à l'état final.}
\objectif{Déterminer un taux d'avancement final et le relier au caractère total ou non total de la transformation.}
\objectif{Exprimer un quotient de réaction et déterminer sa valeur à l'état final.}
\objectif{Montrer que le quotient de réaction à l'équilibre ne dépend pas de la composition initiale.}
\objectif{Déterminer le sens d'évolution spontanée d'un système chimique.}
\prerequis{Tableau d'avancement, avancement maximal, réactif limitant.}
\prerequis{pH, conductivité, absorbance (chapitres~1 et~2).}
\prerequis{Vitesses de réaction, actes élémentaires (chapitres~4 et~5).}
\end{ouverture}
\essentiel
\essSection{1}{Transformations totales et non totales}
\begin{definition}[Taux d'avancement final]
Le taux d'avancement final d'une transformation est le quotient de l'avancement final par l'avancement maximal :
\[\tau=\frac{x_\mathrm{f}}{x_{\max}},\qquad 0\leq\tau\leq1.\]
Si $\tau=1$, la transformation est \textbf{totale} : un réactif a disparu. Si $\tau<1$, elle est \textbf{non totale} :
dans l'état final, tous les réactifs et tous les produits coexistent.
\end{definition}
\begin{exemple}
On dissout $\qty{1,0e-3}{mol}$ d'acide éthanoïque dans 100,0~mL d'eau. L'acide réagit avec l'eau selon
\[\ce{CH3COOH + H2O <=> CH3COO- + H3O+}.\]
Le pH final vaut 3,4, donc $[\ce{H3O+}]_\mathrm{f}=10^{-3{,}4}=\qty{4,0e-4}{mol.L^{-1}}$ et $x_\mathrm{f}=\qty{4,0e-5}{mol}$,
alors que $x_{\max}=\qty{1,0e-3}{mol}$ : $\tau=0{,}040$. Seules 4~\% des molécules d'acide ont réagi avec l'eau.
\end{exemple}
\begin{savoirfaire}[Mettre en évidence le caractère non total]
Ajouter à l'état final un peu de l'un des réactifs : si le système évolue de nouveau (changement de couleur, de pH,
de conductivité), c'est que l'autre réactif était encore présent. La transformation n'était donc pas totale.
\end{savoirfaire}
\essSection{2}{L'équilibre dynamique}
\begin{definition}[État d'équilibre dynamique]
Une transformation non totale est modélisée par deux réactions opposées qui se produisent simultanément ; son
équation s'écrit avec le symbole $\rightleftharpoons$. Dans l'état final, les deux réactions ont la même vitesse : la composition
du système n'évolue plus à l'échelle macroscopique, alors qu'à l'échelle microscopique les deux réactions se
poursuivent. Cet état est appelé \textbf{état d'équilibre dynamique}.
\end{definition}
\schema{07-equilibre-dynamique}
\essSection{3}{Quotient de réaction et constante d'équilibre}
\begin{definition}[Quotient de réaction]
Pour une réaction d'équation $\alpha\,\mathrm{A}+\beta\,\mathrm{B}\rightleftharpoons\gamma\,\mathrm{C}+\delta\,\mathrm{D}$
entre espèces dissoutes, le quotient de réaction est
\[Q_\mathrm{r}=\frac{[\mathrm{C}]^{\gamma}\,[\mathrm{D}]^{\delta}}{[\mathrm{A}]^{\alpha}\,[\mathrm{B}]^{\beta}},\]
où chaque concentration, en \unit{mol.L^{-1}}, est divisée par la concentration standard $c^\circ=\qty{1}{mol.L^{-1}}$ :
$Q_\mathrm{r}$ est sans unité. L'eau, solvant, et les solides n'apparaissent pas dans son expression.
\end{definition}
\begin{propriete}[Constante d'équilibre]
Dans l'état d'équilibre, le quotient de réaction prend une valeur $Q_{\mathrm{r},\text{éq}}=K(T)$, appelée
\textbf{constante d'équilibre}. Elle ne dépend que de la réaction et de la température, et non de la composition
initiale du système. Lorsque $K$ est très grande (en pratique $K>10^4$), un réactif disparaît, ou presque, avant que
$Q_\mathrm{r}$ n'atteigne $K$ : la transformation peut être considérée comme totale.
\end{propriete}
\begin{exemple}
Pour l'acide éthanoïque de l'exemple précédent, $[\ce{CH3COO-}]_\text{éq}=[\ce{H3O+}]_\text{éq}=\qty{4,0e-4}{mol.L^{-1}}$
et $[\ce{CH3COOH}]_\text{éq}=\qty{9,6e-3}{mol.L^{-1}}$, d'où
\[K=\frac{[\ce{CH3COO-}]_\text{éq}\,[\ce{H3O+}]_\text{éq}}{[\ce{CH3COOH}]_\text{éq}}=\frac{(\num{4.0e-4})^2}{\num{9.6e-3}}\approx\num{1.7e-5}\quad\text{à 25~°C}.\]
\end{exemple}
\begin{experience}[Montrer que $K$ ne dépend pas de l'état initial]
On prépare, par dilution d'une solution mère, des solutions d'acide éthanoïque de concentrations apportées $c$
comprises entre \qty{1,0e-3}{mol.L^{-1}} et \qty{1,0e-2}{mol.L^{-1}}. Après étalonnage du conductimètre, on mesure la
conductivité $\sigma$ de chaque solution à 25~°C, de la plus diluée à la plus concentrée, en rinçant la cellule avec
la solution suivante. Les ions \ce{H3O+} et \ce{CH3COO-}, seuls présents en quantité notable, ont la même
concentration ; la loi de Kohlrausch donne alors
\[[\ce{H3O+}]_\text{éq}=\frac{\sigma}{\lambda_{\ce{H3O+}}+\lambda_{\ce{CH3COO-}}},\]
d'où $\tau$ et $Q_{\mathrm{r},\text{éq}}$ pour chaque solution. Le taux d'avancement final augmente avec la dilution, mais
$Q_{\mathrm{r},\text{éq}}$ garde la même valeur, voisine de \num{1.7e-5} : c'est la constante d'équilibre (exercice~7).
\end{experience}
\begin{savoirfaire}[Déterminer la composition d'un système à l'équilibre]
Dresser un tableau d'avancement, exprimer les concentrations à l'équilibre en fonction de l'avancement, écrire
$Q_{\mathrm{r},\text{éq}}=K$, puis résoudre l'équation obtenue : à la main pour une équation du second degré, ou par
dichotomie à l'aide d'un programme.
\end{savoirfaire}
\begin{codecote}[0.56]
\begin{pythoncol}
K, c = 1.7e-5, 1.0e-2   # c en mol/L
def f(x):               # Qr(x) - K
    return x**2 / (c - x) - K
a, b = 0, 0.999 * c     # f(a) < 0 < f(b)
while b - a > 1e-9:
    m = (a + b) / 2
    if f(m) < 0:
        a = m
    else:
        b = m
print(f"x = {a:.2e} ; tau = {a / c:.3f}")
\end{pythoncol}
\begin{consolecol}
x = 4.04e-04 ; tau = 0.040
\end{consolecol}
\cote
\fontsize{9.5}{12.5}\selectfont
\capacite{Capacité numérique}\par\smallskip
Pour l'acide éthanoïque à \qty{1,0e-2}{mol.L^{-1}}, l'avancement volumique à l'équilibre $x$ vérifie
$x^2/(c-x)=K$. La fonction $f$ est croissante sur $[0\,;\,c[$ et change de signe : la dichotomie divise à chaque
étape par deux l'intervalle qui contient la solution. On retrouve un taux d'avancement final de 4~\%.
\end{codecote}
\essSection{4}{Critère d'évolution spontanée}
\begin{theoreme}[Critère d'évolution spontanée]
Un système hors d'équilibre, de quotient de réaction initial $Q_{\mathrm{r},\mathrm{i}}$, évolue spontanément de sorte que
$Q_\mathrm{r}$ tende vers $K$ :
\begin{itemize}[leftmargin=1.4em,itemsep=0.1em]
\item si $Q_{\mathrm{r},\mathrm{i}}<K$, le système évolue dans le \textbf{sens direct} ;
\item si $Q_{\mathrm{r},\mathrm{i}}>K$, le système évolue dans le \textbf{sens inverse} ;
\item si $Q_{\mathrm{r},\mathrm{i}}=K$, le système est à l'équilibre : il n'évolue pas.
\end{itemize}
\end{theoreme}
\schema{07-critere}
\begin{savoirfaire}[Prévoir le sens d'évolution spontanée]
Écrire l'expression de $Q_\mathrm{r}$, calculer $Q_{\mathrm{r},\mathrm{i}}$ avec les concentrations de l'état initial, puis le
comparer à $K$. Si un produit dissous est absent de l'état initial, $Q_{\mathrm{r},\mathrm{i}}=0<K$ : le système évolue dans le
sens direct. Si un réactif est absent, la réaction ne peut se produire que dans le sens inverse.
\end{savoirfaire}
\begin{exemple}
On plonge un fil d'argent dans une solution où $[\ce{Fe^3+}]=\qty{0,10}{mol.L^{-1}}$, $[\ce{Fe^2+}]=\qty{0,050}{mol.L^{-1}}$
et $[\ce{Ag+}]=\qty{0,20}{mol.L^{-1}}$. Pour \ce{Fe^3+(aq) + Ag(s) <=> Fe^2+(aq) + Ag+(aq)}, de constante $K=0{,}31$ à 25~°C,
$Q_{\mathrm{r},\mathrm{i}}=(0{,}050\times0{,}20)/0{,}10=0{,}10<K$ : le système évolue dans le sens direct, l'argent est oxydé.
\end{exemple}
\begin{attention}
La seule valeur de $K$ ne permet pas de prévoir le sens d'évolution : il faut comparer $Q_{\mathrm{r},\mathrm{i}}$ à $K$.
Une même réaction peut évoluer dans un sens ou dans l'autre selon l'état initial du système.
\end{attention}
\begin{rediger}
Écrire l'expression de $Q_\mathrm{r}$ avant de la calculer, préciser l'état (initial ou d'équilibre) des concentrations
utilisées, puis conclure par une phrase complète : « $Q_{\mathrm{r},\mathrm{i}}<K$ : le système évolue spontanément dans le
sens direct de l'équation ».
\end{rediger}
\exercices

\rubrique{Découverte}
\begin{blocLivre}{ex-7-1}
\begin{exercice}[1]{Total ou non total ?}
On donne, pour trois transformations, la composition de l'état final.
\begin{questions}[label=\otLabel{\Alph*.},leftmargin=1.9em,itemsep=0pt,parsep=0.15em]
\item \mbox{\ce{HCl + H2O -> Cl- + H3O+}} : on dissout \qty{1,0e-2}{mol} de chlorure d'hydrogène dans 1,0~L d'eau ;
le pH final vaut 2,0.
\item \mbox{\ce{CH3COOH + H2O <=> CH3COO- + H3O+}} : on dissout \qty{1,0e-2}{mol} d'acide éthanoïque dans 1,0~L
d'eau ; le pH final vaut 3,4.
\item \mbox{\ce{CH3COOH + C2H5OH <=> CH3COOC2H5 + H2O}} : on mélange 0,50~mol d'acide éthanoïque et 0,50~mol
d'éthanol ; il se forme 0,33~mol d'éthanoate d'éthyle.
\end{questions}
\begin{questions}
\item Définir le taux d'avancement final d'une transformation.
\item Calculer le taux d'avancement final de chacune des trois transformations.
\item Lesquelles sont totales ? Lesquelles sont non totales ?
\item Quelles espèces le système B contient-il dans son état final ?
\end{questions}
\end{exercice}
\end{blocLivre}
\begin{blocLivre}{ex-7-2}
\begin{exercice}[1]{Écrire un quotient de réaction}
On considère les réactions suivantes.
\begin{questions}[label=\otLabel{\Alph*.},leftmargin=1.9em,itemsep=0pt,parsep=0.15em]
\item \ce{CH3COOH(aq) + H2O(l) <=> CH3COO^-(aq) + H3O^+(aq)}
\item \ce{Cu^2+(aq) + Zn(s) <=> Cu(s) + Zn^2+(aq)}
\item \ce{AgCl(s) <=> Ag^+(aq) + Cl^-(aq)}
\item \ce{I2(aq) + 2 S2O3^2-(aq) <=> 2 I^-(aq) + S4O6^2-(aq)}
\item \ce{Fe^3+(aq) + SCN^-(aq) <=> FeSCN^2+(aq)}
\end{questions}
\begin{questions}
\item Écrire l'expression du quotient de réaction de chacune de ces réactions.
\item Pourquoi l'eau et les solides n'apparaissent-ils pas dans ces expressions ?
\item Pour la réaction E, calculer $Q_\mathrm{r}$ lorsque
\[[\ce{Fe^3+}]=[\ce{SCN-}]=\qty{2,0e-3}{mol.L^{-1}}\quad\text{et}\quad[\ce{FeSCN^2+}]=\qty{4,0e-4}{mol.L^{-1}}.\]
\end{questions}
\end{exercice}
\end{blocLivre}
\begin{blocLivre}{ex-7-3}
\begin{exercice}[1]{Prévoir le sens d'évolution}
Pour chacun des systèmes suivants, toutes les espèces citées sont présentes dans l'état initial.
\begin{center}\small\renewcommand{\arraystretch}{1.25}
\begin{tabularx}{\linewidth}{@{}lYl@{}}\toprule
\otTete{Réaction} & \otTete{Composition initiale (mol$\cdot$L$^{-1}$)} & \otTete{$K$ à 25~°C}\\\midrule
A. \ce{CH3COOH + H2O <=> CH3COO- + H3O+} & $[\ce{CH3COOH}]=[\ce{CH3COO-}]=0{,}10$ ; $[\ce{H3O+}]=\num{1.0e-4}$ & $\num{1.7e-5}$\\
B. \ce{Fe^3+ + SCN- <=> FeSCN^2+} & toutes les concentrations valent $\num{1.0e-3}$ & $\num{1.0e2}$\\
C. \ce{Cu^2+ + Zn(s) <=> Cu(s) + Zn^2+} & $[\ce{Cu^2+}]=[\ce{Zn^2+}]=0{,}10$ & $\num{1.9e37}$\\
D. \ce{AgCl(s) <=> Ag+ + Cl-} & $[\ce{Ag+}]=[\ce{Cl-}]=\num{1.0e-3}$ (sans solide) & $\num{1.8e-10}$\\\bottomrule
\end{tabularx}\end{center}
\begin{questions}
\item Calculer le quotient de réaction initial de chaque système.
\item Prévoir le sens d'évolution spontanée de chaque système.
\item Qu'observe-t-on lorsqu'on prépare le mélange D ?
\end{questions}
\end{exercice}
\end{blocLivre}
\begin{blocLivre}{ex-7-4}
\begin{exercice}[1]{Ajouter un réactif à l'équilibre}
On mélange des solutions très diluées de chlorure de fer(III) et de thiocyanate de potassium. La solution prend une
teinte rouge orangé, due aux ions \ce{FeSCN^2+} formés selon la réaction d'équation
\[\ce{Fe^3+(aq) + SCN^-(aq) <=> FeSCN^2+(aq)}.\]
Le mélange, à l'équilibre, est réparti dans trois tubes ; on ajoute quelques gouttes d'une solution concentrée de
chlorure de fer(III) dans le tube 2, de thiocyanate de potassium dans le tube 3.
\schema{07-ex04}
\begin{questions}
\item Comment la couleur évolue-t-elle dans les tubes 2 et 3 ? Qu'en déduire pour la quantité d'ions \ce{FeSCN^2+} ?
\item Quel réactif était encore présent dans le mélange avant l'ajout d'ions \ce{Fe^3+} ? avant l'ajout d'ions \ce{SCN-} ?
\item La transformation qui a conduit au mélange initial était-elle totale ? Justifier.
\item Justifier l'évolution observée dans le tube 2 à l'aide du quotient de réaction.
\end{questions}
\end{exercice}
\end{blocLivre}
\begin{blocLivre}{ex-7-5}
\begin{exercice}[1]{Un équilibre dynamique}
Le graphique représente les vitesses des deux réactions opposées qui modélisent la transformation
\ce{A + B <=> C + D}, à partir d'un mélange initial de A et de B.
\schema{07-ex05}
\begin{questions}
\item Laquelle des deux courbes représente la vitesse de la réaction dans le sens direct ? Justifier.
\item À partir de quelle date le système peut-il être considéré à l'équilibre ?
\item Que se passe-t-il à l'échelle microscopique dans l'état d'équilibre ? Pourquoi parle-t-on d'équilibre dynamique ?
\item La composition du système varie-t-elle après 200~min ? Justifier.
\end{questions}
\end{exercice}
\end{blocLivre}
\rubrique{Entraînement}
\begin{blocLivre}{ex-7-6}
\begin{exercice}[2]{Le taux d'avancement dépend de l'état initial}
On prépare deux solutions d'acide méthanoïque \ce{HCOOH} : $S_1$ à $c_1=\qty{1,0e-2}{mol.L^{-1}}$ et $S_2$ à
$c_2=\qty{1,0e-3}{mol.L^{-1}}$. Leurs pH valent respectivement 2,90 et 3,46.
\begin{questions}
\item Écrire l'équation de la réaction de l'acide méthanoïque avec l'eau et dresser le tableau d'avancement pour un
volume $V$ de solution.
\item Calculer le taux d'avancement final de la réaction dans chacune des solutions.
\item Calculer le quotient de réaction à l'équilibre dans chacune des solutions.
\item Que peut-on conclure de la comparaison des résultats des questions 2 et 3 ?
\end{questions}
\end{exercice}
\end{blocLivre}
\begin{blocLivre}{ex-7-7}
\begin{exercice}[2]{Une constante indépendante de l'état initial}
On mesure, à 25~°C, la conductivité de solutions d'acide éthanoïque de différentes concentrations $c$.
\begin{center}\small\renewcommand{\arraystretch}{1.2}
\begin{tabular}{@{}l*{4}{c}@{}}\toprule
$c$ ($10^{-3}$ mol$\cdot$L$^{-1}$) & 10,0 & 5,0 & 2,0 & 1,0\\
$\sigma$ ($\upmu$S$\cdot$cm$^{-1}$) & 160 & 112 & 69,6 & 48,3\\\bottomrule
\end{tabular}\end{center}
\begin{donnees}
$\lambda(\ce{H3O+})=\qty{35,0}{mS.m^2.mol^{-1}}$ ; $\lambda(\ce{CH3COO-})=\qty{4,09}{mS.m^2.mol^{-1}}$ ; la concentration
en ions hydroxyde est négligeable ; 1~$\upmu$S$\cdot$cm$^{-1}$ $=10^{-4}$ S$\cdot$m$^{-1}$.
\end{donnees}
\begin{questions}
\item Montrer que $\sigma=\left(\lambda_{\ce{H3O+}}+\lambda_{\ce{CH3COO-}}\right)[\ce{H3O+}]_\text{éq}$.
\item Calculer $[\ce{H3O+}]_\text{éq}$, en \unit{mol.L^{-1}}, pour chaque solution.
\item Calculer le quotient de réaction à l'équilibre pour chaque solution. Conclure et donner la valeur de $K$.
\item Calculer le taux d'avancement final pour chaque solution. Comment évolue-t-il avec la dilution ?
\end{questions}
\end{exercice}
\end{blocLivre}
\begin{blocLivre}{ex-7-8}
\begin{exercice}[2]{Deux sens d'évolution pour une même réaction}
On étudie la réaction \ce{Fe^3+(aq) + Ag(s) <=> Fe^2+(aq) + Ag^+(aq)}, de constante $K=0{,}31$ à 25~°C. On prépare deux
systèmes contenant chacun de la poudre d'argent en excès.
\begin{center}\small\renewcommand{\arraystretch}{1.2}
\begin{tabular}{@{}l*{3}{c}@{}}\toprule
 & $[\ce{Fe^3+}]_\mathrm{i}$ (mol$\cdot$L$^{-1}$) & $[\ce{Fe^2+}]_\mathrm{i}$ (mol$\cdot$L$^{-1}$) & $[\ce{Ag+}]_\mathrm{i}$ (mol$\cdot$L$^{-1}$)\\\midrule
système A & 0,10 & 0,010 & 0,010\\
système B & 0,010 & 0,10 & 0,50\\\bottomrule
\end{tabular}\end{center}
\begin{questions}
\item Écrire l'expression du quotient de réaction.
\item Déterminer le sens d'évolution de chaque système. Qu'observe-t-on sur la poudre d'argent dans chaque cas ?
\item Pour le système A, montrer que l'avancement volumique à l'équilibre $x$ vérifie
\[\frac{(0{,}010+x)^2}{0{,}10-x}=0{,}31,\]
puis le calculer. En déduire $[\ce{Fe^3+}]_\text{éq}$ et le taux d'avancement final.
\item La constante d'équilibre est inférieure à 1, et pourtant le taux d'avancement final du système A dépasse 75~\%.
Commenter.
\end{questions}
\end{exercice}
\end{blocLivre}
\begin{blocLivre}{ex-7-9}
\begin{exercice}[2]{Le diiode et les ions iodure}
Le diiode est peu soluble dans l'eau, mais il réagit avec les ions iodure pour former des ions triiodure :
\ce{I2(aq) + I^-(aq) <=> I3^-(aq)}, de constante $K=\num{7.1e2}$ à 25~°C. C'est ainsi que l'on prépare des solutions
iodées concentrées. Le programme suivant détermine la composition d'un mélange à l'équilibre.
\begin{python}
K = 7.1e2
c1, c2 = 1.0e-3, 1.0e-2    # [I2] et [I-] initiales, en mol/L
def f(x):
    return ...             # Qr(x) - K
a, b = 0, 0.999 * c1
while b - a > 1e-10:
    m = (a + b) / 2
    if f(m) < 0:
        a = m
    else:
        b = m
print(f"x = {a:.3e} ; tau = {a / c1:.3f}")
\end{python}
L'exécution affiche la première ligne ci-dessous ; en remplaçant \code{c2} par \code{1.0e-1}, on obtient la seconde.
\begin{console}
x = 8.664e-04 ; tau = 0.866
x = 9.860e-04 ; tau = 0.986
\end{console}
\begin{questions}
\item \capacite{Capacité numérique}Dresser le tableau d'avancement et exprimer $Q_\mathrm{r}$ en fonction de l'avancement volumique $x$.
\item Compléter la ligne incomplète du programme.
\item Expliquer le principe de la boucle \code{while}. Pourquoi faut-il que $f(a)$ et $f(b)$ soient de signes contraires ?
\item Interpréter les deux résultats affichés.
\end{questions}
\end{exercice}
\end{blocLivre}
\begin{blocLivre}{ex-7-10}
\begin{exercice}[2]{L'estérification, une transformation limitée}
L'acide éthanoïque réagit avec l'éthanol selon la réaction d'équation
\[\ce{CH3COOH(l) + C2H5OH(l) <=> CH3COOC2H5(l) + H2O(l)}.\]
Toutes les espèces forment un mélange homogène, dans lequel on écrit le quotient de réaction avec les concentrations,
y compris celle de l'eau ; la constante d'équilibre vaut $K=4{,}0$.
\begin{questions}
\item Montrer que le quotient de réaction peut s'écrire avec les quantités de matière :
\[Q_\mathrm{r}=\frac{n(\text{ester})\times n(\text{eau})}{n(\text{acide})\times n(\text{alcool})}.\]
\item On mélange 1,00~mol d'acide et 1,00~mol d'alcool. Déterminer l'avancement à l'équilibre et le taux
d'avancement final.
\item On mélange 1,00~mol d'acide et 2,00~mol d'alcool. Montrer que l'avancement à l'équilibre vérifie
$3x^2-12x+8=0$, puis calculer le taux d'avancement final.
\item Au laboratoire, on élimine parfois l'eau au fur et à mesure de sa formation. Quel est l'intérêt de cette
technique ? Justifier à l'aide du critère d'évolution.
\end{questions}
\end{exercice}
\end{blocLivre}
\begin{blocLivre}{ex-7-11}
\begin{exercice}[2]{Une transformation totale}
On plonge du zinc en poudre dans 100~mL d'une solution de sulfate de cuivre(II) à $\qty{0,10}{mol.L^{-1}}$. La réaction
\ce{Cu^2+(aq) + Zn(s) <=> Cu(s) + Zn^2+(aq)} a pour constante $K=\num{1.9e37}$ à 25~°C.
\begin{donnees}
$M(\ce{Zn})=\qty{65,4}{g.mol^{-1}}$ ; $N_\mathrm{A}=\qty{6,02e23}{mol^{-1}}$.
\end{donnees}
\begin{questions}
\item Écrire l'expression du quotient de réaction.
\item Le zinc est introduit en large excès : en fin de transformation, $[\ce{Zn^2+}]=\qty{0,10}{mol.L^{-1}}$. Calculer
$[\ce{Cu^2+}]$ à l'équilibre, puis le nombre d'ions \ce{Cu^2+} restant dans la solution. Conclure.
\item On introduit seulement 0,50~g de zinc. Quel est le réactif limitant ? Calculer le quotient de réaction en fin de
transformation et le comparer à $K$. L'état d'équilibre est-il atteint ? Pourquoi la transformation s'arrête-t-elle ?
\item Justifier qu'une transformation de constante $K>10^4$ peut être considérée comme totale.
\end{questions}
\end{exercice}
\end{blocLivre}
\rubrique{Objectif bac}
\begin{blocLivre}{ex-7-12}
\begin{exercice}[3]{Mesurer une constante d'équilibre}
\begin{docu}{1}{Trois mélanges}
On réalise trois mélanges de solutions de chlorure de fer(III) et de thiocyanate de potassium. Les concentrations
apportées dans chaque mélange, avant réaction, sont données ci-dessous. Une fois l'équilibre
\ce{Fe^3+(aq) + SCN^-(aq) <=> FeSCN^2+(aq)} atteint, on mesure l'absorbance à 470~nm, dans une cuve de 1,00~cm.
\par\smallskip
{\small\renewcommand{\arraystretch}{1.2}
\begin{tabular}{@{}l*{3}{c}@{}}\toprule
 & mélange 1 & mélange 2 & mélange 3\\\midrule
$[\ce{Fe^3+}]_\mathrm{i}$ ($10^{-3}$ mol$\cdot$L$^{-1}$) & 1,0 & 2,0 & 1,0\\
$[\ce{SCN-}]_\mathrm{i}$ ($10^{-3}$ mol$\cdot$L$^{-1}$) & 1,0 & 1,0 & 3,0\\
$A$ à 470~nm & 0,394 & 0,732 & 1,02\\\bottomrule
\end{tabular}}
\end{docu}
\begin{donnees}
À 470~nm, seuls les ions \ce{FeSCN^2+} absorbent ; leur coefficient d'absorption molaire vaut
\mbox{$\varepsilon=\qty{4,7e3}{L.mol^{-1}.cm^{-1}}$}.
\end{donnees}
\begin{questions}
\item Pourquoi effectue-t-on les mesures à 470~nm ? Rappeler la loi de Beer-Lambert.
\item Pour le mélange 1, déterminer $[\ce{FeSCN^2+}]_\text{éq}$, puis $[\ce{Fe^3+}]_\text{éq}$ et $[\ce{SCN-}]_\text{éq}$.
\item \initiative Montrer que les trois mélanges conduisent à la même valeur du quotient de réaction à l'équilibre.
En déduire la constante d'équilibre.
\item Calculer le taux d'avancement final pour les mélanges 1 et 3. Comment un excès d'ions thiocyanate agit-il sur
la formation des ions \ce{FeSCN^2+} ?
\end{questions}
\end{exercice}
\end{blocLivre}
\begin{blocLivre}{ex-7-13}
\begin{exercice}[3]{Dissoudre un précipité de chlorure d'argent}
\begin{docu}{1}{Un solide peu soluble}
Le chlorure d'argent est un solide blanc très peu soluble dans l'eau : \ce{AgCl(s) <=> Ag^+(aq) + Cl^-(aq)},
de constante $K_1=\num{1.8e-10}$ à 25~°C. Il se dissout pourtant dans une solution d'ammoniac, en formant des ions
diammineargent(I) : \ce{AgCl(s) + 2 NH3(aq) <=> Ag(NH3)2^+(aq) + Cl^-(aq)}, de constante $K_2=\num{2.9e-3}$.
\end{docu}
\begin{donnees}
$M(\ce{AgCl})=\qty{143,3}{g.mol^{-1}}$.
\end{donnees}
\begin{questions}
\item Écrire l'expression du quotient de réaction de chacune des deux réactions.
\item On appelle solubilité la concentration maximale du solide dissous. Calculer la solubilité du chlorure d'argent
dans l'eau pure, en \unit{mol.L^{-1}} puis en \unit{g.L^{-1}}.
\item On ajoute du chlorure d'argent à une solution d'ammoniac à $\qty{1,0}{mol.L^{-1}}$. Justifier que le solide
commence à se dissoudre.
\item \initiative Déterminer la masse maximale de chlorure d'argent qui peut se dissoudre dans 1,0~L de cette
solution. Comparer avec l'eau pure.
\end{questions}
\end{exercice}
\end{blocLivre}
\begin{blocLivre}{ex-7-14}
\begin{exercice}[3]{Extraire le diiode d'une solution aqueuse}
\begin{docu}{1}{Une extraction liquide-liquide}
Pour récupérer le diiode d'une solution aqueuse, on l'agite avec du cyclohexane, solvant non miscible à l'eau, dans
une ampoule à décanter. Le diiode se répartit entre les deux phases selon l'équilibre
\ce{I2(aq) <=> I2(cyclo)}, de constante \mbox{$K=[\ce{I2}]_\text{cyclo}/[\ce{I2}]_\text{aq}=50$} à 25~°C.
\end{docu}
\schema{07-ex14}
\begin{donnees}
On traite 100~mL d'une solution aqueuse de diiode à $\qty{1,0e-3}{mol.L^{-1}}$ ; on dispose de 20~mL de cyclohexane.
\end{donnees}
\begin{questions}
\item Pourquoi la phase organique se trouve-t-elle au-dessus de la phase aqueuse ?
\item Montrer qu'après une extraction avec un volume $V_\mathrm{o}$ de cyclohexane, la fraction du diiode restant dans
la phase aqueuse, de volume $V_\mathrm{a}$, vaut $\dfrac{1}{1+K\,V_\mathrm{o}/V_\mathrm{a}}$. Calculer cette fraction et la
concentration du diiode dans chaque phase pour une extraction avec 20~mL de cyclohexane.
\item \initiative Vaut-il mieux réaliser une extraction avec 20~mL de cyclohexane ou deux extractions successives avec
10~mL ? Justifier par le calcul.
\item Proposer une explication à la grande solubilité du diiode dans le cyclohexane.
\end{questions}
\end{exercice}
\end{blocLivre}
\begin{blocLivre}{ex-7-15}
\begin{exercice}[3]{La mutarotation du glucose}
\begin{docu}{1}{Deux formes du glucose}
En solution aqueuse, le glucose existe principalement sous deux formes cycliques, notées $\alpha$ et $\beta$, qui se
transforment lentement l'une en l'autre : glucose $\alpha$ \ce{<=>} glucose $\beta$. On modélise cette transformation
par deux réactions opposées de vitesses $v_1=k_1[\alpha]$ et $v_{-1}=k_{-1}[\beta]$, avec $k_1=\qty{1,6e-2}{min^{-1}}$ et
$k_{-1}=\qty{9,0e-3}{min^{-1}}$ (valeurs adoptées à 20~°C).
\end{docu}
\begin{docu}{2}{Simulation}
Le programme suivant calcule l'évolution des proportions $a$ de forme $\alpha$ et $b$ de forme $\beta$, à partir de
glucose $\alpha$ pur, par pas de 1~min.
\begin{python}
k1, k2, dt = 0.016, 0.009, 1.0
a, b = 1.0, 0.0
for i in range(240):
    v1, v2 = k1 * a, k2 * b
    a = a + (v2 - v1) * dt
    b = b + (v1 - v2) * dt
print(round(a, 3), round(b, 3), round(b / a, 2))
\end{python}
L'exécution affiche \code{0.361 0.639 1.77}. Le graphique compare cette simulation à une simulation partant de
glucose $\beta$ pur.
\end{docu}
\schema{07-ex15}
\begin{questions}
\item Exprimer le quotient de réaction de la transformation en fonction des proportions $a$ et $b$.
\item \capacite{Capacité numérique}Expliquer les deux lignes qui mettent à jour \code{a} et \code{b}. Pourquoi la somme
$a+b$ reste-t-elle égale à 1 ?
\item Quelle est la composition du mélange à l'équilibre ? Dépend-elle de l'état initial ?
\item Montrer qu'à l'équilibre $K=k_1/k_{-1}$. Calculer $K$ et comparer à la valeur affichée par le programme.
\item \initiative Un catalyseur multiplie $k_1$ et $k_{-1}$ par le même facteur. Quel est son effet sur la composition
finale ? sur la durée nécessaire pour l'atteindre ?
\end{questions}
\end{exercice}
\end{blocLivre}
